\section{A primitive mixed coordinate in degree eleven}
\label{sec:eleven-mixed-mixed}

We remove a primitive mixed coefficient while preserving the arithmetic
factorization. Use the groups $G$, $H$, $N$, $D$, and $D_{10}$ defined above.
Write $M=\ker\lambda$ and
\[
 C_\infty=\{g\in D:B_g(X,Y)=1\},\qquad
 E=C_\infty\cap D_{10},\qquad J(E)=E\cap H.
\]
Here $B_g$ is the dyadic beta multiplier. The preceding normalizations
give $G=EH$, and every Magnus series $f_g$ with $g\in E$ has no
nonconstant term below degree eleven. We use these established facts
and the exact degree-eight leading lattice below. The argument uses one coordinate of the degree-eleven leading space.

\begin{theorem}\label{thm:eleven-mixed-mixed-reduction}
The coordinate
\[
 b_{11}(g)=[S^6TST^3]f_g\qquad(g\in E)
\]
is a continuous additive homomorphism satisfying
\[
 b_{11}(aga^{-1})=\lambda_2(a)^{11}b_{11}(g)
 \quad(a\in G,\ g\in E).
\]
Its full and arithmetic images are both saturated:
\[
 b_{11}(E)=b_{11}(J(E))=\mathbf Z_2.
\]
Consequently
\[
 E'=\ker(b_{11}|_E)
\]
is a proper closed normal subgroup of $G$, and
\[
 E/E'\simeq\mathbf Z_2,\qquad G=E'H,
 \qquad
 B\simeq E'/\overline{[G,E'](E'\cap H)}.
\]
For $X_2=H^1(N,\mathbf F_2)^G$, its restriction kernel remains
\[
 \ker\bigl(X_2\longrightarrow H^1(E',\mathbf F_2)\bigr)
 =\left(\bigoplus_{q\text{ odd prime}}\mathbf F_2\psi_q\right)
       \oplus\mathbf F_2\eta_3.
\]
The invariant character $b_{11}\bmod2$ on $E$ is nonzero on $J(E)$
and has no $G$-invariant extension to $N$.
\end{theorem}

\begin{proof}
We first give the integral coefficient calculation that supplies an
arithmetic unit. For Lie polynomials use
\[
 P_3=[S,[S,T]]+[T,[S,T]],\qquad
 \{A,B\}=[A,B]+\partial_A B-\partial_B A,
\]
where $\partial_A(S)=0$ and $\partial_A(T)=[T,A]$.
This is the Ihara bracket in the convention of the GT product.
Let $L_8$ be the integer Lie polynomial for the already established
degree-eight leading line, normalized so that
\[
 [S^4TST^2]L_8=7.
\]
For specificity it is completely specified by its coordinates in the
lexicographically ordered Lyndon basis of weight eight, with $S<T$:
\begin{align*}
(&0,0,-4,0,-2,7,17,0,-1,-12,-6,4,-3,-6,0,\\
 &-5,3,3,6,7,21,3,-15,-3,0,8,-11,-4,-10,0).
\end{align*}
The Lyndon bracket of a nonletter word is $[L_u,L_v]$, where $v$
is its longest proper Lyndon suffix and the word is $uv$.
This finite recursion gives, in particular,
\begin{equation}\label{eq:eleven-mixed-four-coefficients}
 [S^4TST^2]L_8=7,\qquad [S^5TST]L_8=-4,\qquad
 [S^5T^3]L_8=[S^6T^2]L_8=0.
\end{equation}
An element $z\in D_7$ therefore has degree-eight term
$(m_8(z)/7)L_8$.

Put $\omega=S^6TST^3$. Expanding
\[
 P_3=S^2T-2STS+TS^2+2TST-T^2S-ST^2
\]
shows that $[\omega]P_3L_8=[\omega]L_8P_3=0$: the required
prefix or suffix of length three is absent from $P_3$. Moreover
\[
 [T,P_3]=-2TSTS+T^2S^2+3T^2ST-T^3S-3TST^2
             -S^2T^2+2STST+ST^3.
\]
Replacing one $T$ in a word of $L_8$ by this expression gives
\[
 [\omega]\partial_{P_3}L_8
     =2[S^5T^3]L_8-2[S^6T^2]L_8=0.
\]
For $\partial_{L_8}P_3$, the only surviving contributions are
$-S^2L_8T$ from $S^2T$ and $+SL_8T^2$ from $-ST^2$.
Thus the final, negative derivation term contributes
\[
 -[\omega]\partial_{L_8}P_3
    =[S^4TST^2]L_8-[S^5TST]L_8=7+4=11.
\]
The four Ihara summands have coefficients $(0,0,0,11)$, proving
\begin{equation}\label{eq:eleven-mixed-bracket-unit}
 [S^6TST^3]\{P_3,L_8\}=11.
\end{equation}

Choose the arithmetic elements from the earlier exact
normalizations: $\sigma\in H\cap M$ with
$\rho_2^I(\sigma)=2$ and all odd-prime rational Kummer coordinates
zero, and $h\in H$ with $t=\lambda(h)$ and $t_2=3$.
Let
\[
 a=m_3(\sigma),\qquad u=\tau_5(\sigma).
\]
The established integral moment parity gives $a,u\in\mathbf Z_2^\times$.
The arithmetic word
\[
 w_8=[\sigma,h\sigma h^{-1}]
\]
belongs to $D_7\cap H$ and satisfies the exact earlier formula
\begin{equation}\label{eq:eleven-mixed-old-eight}
 m_8(w_8)=63au.
\end{equation}
Recall that all rational Kummer coordinates and the dyadic beta
multiplier are homomorphisms on $M$, the latter into a commutative
group. In particular $B_{w_8}=1$.

Define the new arithmetic commutator
\[
 w_{11}=[\sigma,w_8].
\]
The first nonconstant terms of $f_\sigma$ and $f_{w_8}$ have
degrees three and eight. The filtered GT commutator has no term
below degree eleven, and its term in degree eleven is their
Ihara bracket. Its rational Kummer coordinates and beta multiplier
are trivial, since both inputs lie in $M$. Hence $w_{11}\in E\cap H$.
Equations \eqref{eq:eleven-mixed-bracket-unit} and
\eqref{eq:eleven-mixed-old-eight} now give the exact value
\begin{equation}\label{eq:eleven-mixed-arithmetic-unit}
 b_{11}(w_{11})
  =\frac{11}{7}a\,m_8(w_8)
  =99a^2u\in\mathbf Z_2^\times.
\end{equation}
There are no contributions from higher terms of either input to
degree eleven.

On a subgroup whose words start in degree eleven, the degree-eleven
part of the GT product is the sum of the two leading parts.
Conjugation by an arbitrary element of $G$ multiplies this
homogeneous part by $\lambda_2^{11}$: on the associated graded
algebra both generators are multiplied by $\lambda_2$, and all
conjugating corrections have higher degree. This proves the
claimed homomorphism and weight laws. Coefficient extraction is
continuous. The closed arithmetic image contains the image of the
procyclic group generated by $w_{11}$, which is all of $\mathbf Z_2$
by \eqref{eq:eleven-mixed-arithmetic-unit}. This proves both image
equalities and the normality and properness of $E'$.

For each $e\in E$, choose $r\in\widehat{\mathbf Z}$ whose
2-adic component is $b_{11}(e)/b_{11}(w_{11})$. The legitimate
profinite power $w_{11}^r$ lies in $J(E)$ and has the same
coordinate as $e$. Thus $e w_{11}^{-r}\in E'$, proving
$E=E'J(E)$ and hence $G=E'H$.

For completeness, let $K=N\cap H$ and
$L=\overline{[G,E']}$. Since $N=E'K$, the image of $E'$ is
central in $G/L$, while $[H,K]\subseteq K$. Consequently
\[
 \overline{[G,N]K}=LK.
\]
The product is closed because both factors are compact and $L$
is normal. Also $E'\cap LK=L(E'\cap H)$. The inclusion of $E'$
therefore induces the stated canonical presentation of $B$.

Finally let $\chi\in X_2$. Since $w_8\in N$ and $\sigma\in G$,
invariance gives
\[
 \chi(w_{11})=\chi(\sigma w_8\sigma^{-1})-\chi(w_8)=0.
\]
If $\chi$ vanishes on $E'$, the same correction by $w_{11}^r$
shows that it vanishes on $E$. The converse is immediate, so the
restriction kernels to $E$ and $E'$ coincide. The previously proved
kernel for $E$ is exactly the displayed rational Kummer span plus
$\mathbf F_2\eta_3$. On the other hand $b_{11}\bmod2$ is invariant
and takes value one on $w_{11}$, proving both its arithmetic
nonvanishing and its nonextension to $N$.
\end{proof}

\begin{remark}\label{rem:eleven-mixed-scope}
This theorem removes a saturated arithmetic coordinate from the
beta-trivial group. It does not produce a nonzero character of the
relative quotient $B$. It also does not assert that the words of
$E'$ start in degree twelve: that stronger conclusion requires a
separate proof that $b_{11}$ detects the whole degree-eleven
beta-trivial leading space.
\end{remark}
